% GENERATED FILE. Edit canonical lessons, then run npm run generate:books.
% canonical-source-hash: fnv1a64:04a513fa97b8473b
% canonical-lessons: IT-C288-annuncio, IT-C288-statistica, IT-C288-rivelare, IT-C288-diffondere, IT-C288-approfondire, IT-R288-ripasso-281-288

\chapter{La notizia si diffonde}
\label{ch:it-la-notizia-si-diffonde}
\hlchaptermodality{288}

\begin{chapteropening}
\textbf{By the end of this chapter:} \emph{I can say that news was announced, revealed and spread, and that someone looked into it more deeply.}

Complete the last lesson of chapter 288: I can say that news was announced, revealed and spread, and that someone looked into it more deeply.
\end{chapteropening}

\section[l'annuncio]{l'annuncio — the announcement; the advert}

\label{lesson:IT-C288-annuncio}



\begin{quote}
Before the new one: say the Italian for the TV channel, then the Italian for the broadcaster.
\end{quote}

\subsection*{You'll want to know: l'annuncio}
\textbf{l'annuncio} — \textquotedblleft{}the announcement; the advert\textquotedblright{}.

\begin{quote}
L'annuncio è arrivato alle otto di sera.
\end{quote}

\begin{quote}
\emph{The announcement came at eight in the evening.}
\end{quote}

In a paper, \emph{gli annunci} are the small ads.

\begin{cousinweb}
\emph{annuncio} is built on \emph{annunciare} \textquotedblleft{}to announce\textquotedblright{}, from Latin \emph{annūntiāre}, built on \emph{nūntius} \textquotedblleft{}messenger\textquotedblright{} (Lewis \& Short).
\end{cousinweb}

\begin{grammarlens}[title={What you've built}]
One more word for telling, explaining and describing at B1.
\end{grammarlens}

\subsection*{Your turn}
\begin{itemize}
\raggedright
  \item \emph{Say it:} \emph{l'annuncio}
  \item \emph{Say it:} \emph{l'annuncio}, once more
  \item \emph{Say it:} a sentence of your own with \emph{l'annuncio}
  \item \emph{Recall:} say the Italian for the TV channel, then the Italian for the broadcaster, then say \emph{l'annuncio} again
\end{itemize}

\begin{tcolorbox}[breakable,colback=teal!4,colframe=teal!35!black,title={Before you move on}]
What does l'annuncio mean? (\textquotedblleft{}The announcement; the advert\textquotedblright{}.) Say it once more.
\end{tcolorbox}
\section[la statistica]{la statistica — the statistic; statistics}

\label{lesson:IT-C288-statistica}



\begin{quote}
Before the new one: say the Italian for the broadcaster, then the Italian for the announcement.
\end{quote}

\subsection*{You'll want to know: la statistica}
\textbf{la statistica} — \textquotedblleft{}the statistic; statistics\textquotedblright{}.

\begin{quote}
Secondo le statistiche, il turismo è cresciuto.
\end{quote}

\begin{quote}
\emph{According to the statistics, tourism has grown.}
\end{quote}

\emph{Secondo le statistiche} — according to the figures.

\begin{cousinweb}
\emph{statistica} is usually traced to Modern Latin \emph{statisticus} \textquotedblleft{}of the state\textquotedblright{}, built on Latin \emph{status} \textquotedblleft{}condition, state\textquotedblright{} (Lewis \& Short): first, figures about a state.
\end{cousinweb}

\begin{grammarlens}[title={What you've built}]
One more word for telling, explaining and describing at B1.
\end{grammarlens}

\subsection*{Your turn}
\begin{itemize}
\raggedright
  \item \emph{Say it:} \emph{la statistica}
  \item \emph{Say it:} \emph{la statistica}, once more
  \item \emph{Say it:} a sentence of your own with \emph{la statistica}
  \item \emph{Recall:} say the Italian for the broadcaster, then the Italian for the announcement, then say \emph{la statistica} again
\end{itemize}

\begin{tcolorbox}[breakable,colback=teal!4,colframe=teal!35!black,title={Before you move on}]
What does la statistica mean? (\textquotedblleft{}The statistic; statistics\textquotedblright{}.) Say it once more.
\end{tcolorbox}
\section[rivelare]{rivelare — to reveal}

\label{lesson:IT-C288-rivelare}



\begin{quote}
Before the new one: say the Italian for the announcement, then the Italian for the statistic.
\end{quote}

\subsection*{You'll want to know: rivelare}
\textbf{rivelare} — \textquotedblleft{}to reveal\textquotedblright{}.

\begin{quote}
Il giornale ha rivelato il nome del vincitore.
\end{quote}

\begin{quote}
\emph{The paper revealed the winner's name.}
\end{quote}

\emph{Rivelarsi} — to turn out: \emph{si è rivelato un errore}.

\begin{cousinweb}
\emph{rivelare} is a learned borrowing from Latin \emph{revēlāre} \textquotedblleft{}to unveil\textquotedblright{} (Lewis \& Short), built on \emph{vēlum} \textquotedblleft{}veil\textquotedblright{}.
\end{cousinweb}

\begin{grammarlens}[title={What you've built}]
One more word for telling, explaining and describing at B1.
\end{grammarlens}

\subsection*{Your turn}
\begin{itemize}
\raggedright
  \item \emph{Say it:} \emph{rivelare}
  \item \emph{Say it:} \emph{rivelare}, once more
  \item \emph{Say it:} a sentence of your own with \emph{rivelare}
  \item \emph{Recall:} say the Italian for the announcement, then the Italian for the statistic, then say \emph{rivelare} again
\end{itemize}

\begin{tcolorbox}[breakable,colback=teal!4,colframe=teal!35!black,title={Before you move on}]
What does rivelare mean? (\textquotedblleft{}To reveal\textquotedblright{}.) Say it once more.
\end{tcolorbox}
\section[diffondere]{diffondere — to spread (news)}

\label{lesson:IT-C288-diffondere}



\begin{quote}
Before the new one: say the Italian for the statistic, then the Italian for to reveal.
\end{quote}

\subsection*{You'll want to know: diffondere}
\textbf{diffondere} — \textquotedblleft{}to spread (news)\textquotedblright{}.

\begin{quote}
La notizia si è diffusa in pochi minuti.
\end{quote}

\begin{quote}
\emph{The news spread in a few minutes.}
\end{quote}

Past participle \emph{diffuso}: \emph{una notizia molto diffusa}.

\begin{cousinweb}
\emph{diffondere} is a learned borrowing from Latin \emph{diffundere} \textquotedblleft{}to pour out, to spread\textquotedblright{} (Lewis \& Short), from \emph{dis-} and \emph{fundere} \textquotedblleft{}to pour\textquotedblright{}.
\end{cousinweb}

\begin{grammarlens}[title={What you've built}]
One more word for telling, explaining and describing at B1.
\end{grammarlens}

\subsection*{Your turn}
\begin{itemize}
\raggedright
  \item \emph{Say it:} \emph{diffondere}
  \item \emph{Say it:} \emph{diffondere}, once more
  \item \emph{Say it:} a sentence of your own with \emph{diffondere}
  \item \emph{Recall:} say the Italian for the statistic, then the Italian for to reveal, then say \emph{diffondere} again
\end{itemize}

\begin{tcolorbox}[breakable,colback=teal!4,colframe=teal!35!black,title={Before you move on}]
What does diffondere mean? (\textquotedblleft{}To spread (news)\textquotedblright{}.) Say it once more.
\end{tcolorbox}
\section[approfondire]{approfondire — to go into in depth, to look more deeply into}

\label{lesson:IT-C288-approfondire}



\begin{quote}
Before the new one: say the Italian for to reveal, then the Italian for to spread.
\end{quote}

\subsection*{You'll want to know: approfondire}
\textbf{approfondire} — \textquotedblleft{}to go into in depth, to look more deeply into\textquotedblright{}.

\begin{quote}
Domani approfondiamo questo argomento.
\end{quote}

\begin{quote}
\emph{Tomorrow we will go into this topic in depth.}
\end{quote}

Built on \emph{profondo}, ‘deep’: to make deeper.

\begin{cousinweb}
\emph{approfondire} is built on \emph{profondo} \textquotedblleft{}deep\textquotedblright{}, from Latin \emph{profundus} (Lewis \& Short), with \emph{ad-} and the verb ending \emph{-ire}.
\end{cousinweb}

\begin{grammarlens}[title={What you've built}]
One more word for telling, explaining and describing at B1.
\end{grammarlens}

\subsection*{Your turn}
\begin{itemize}
\raggedright
  \item \emph{Say it:} \emph{approfondire}
  \item \emph{Say it:} \emph{approfondire}, once more
  \item \emph{Say it:} a sentence of your own with \emph{approfondire}
  \item \emph{Recall:} say the Italian for to reveal, then the Italian for to spread, then say \emph{approfondire} again
\end{itemize}

\begin{tcolorbox}[breakable,colback=teal!4,colframe=teal!35!black,title={Before you move on}]
What does approfondire mean? (\textquotedblleft{}To go into in depth, to look more deeply into\textquotedblright{}.) Say it once more.
\end{tcolorbox}
\section[(dialogue)]{Ripasso dei capitoli 281–288}

\label{lesson:IT-R288-ripasso-281-288}



\begin{quote}
Say the Italian for the eve and the next day.
\end{quote}

\subsection*{Your turn}
Say these aloud:

\begin{itemize}
\raggedright
  \item the eve, the next day, this morning, last night, the whole morning
  \item the news pages, the daily newspaper, the weekly, the headline, the article
  \item the newsroom, the news reporter, the correspondent, the interview, to interview
  \item the inquiry, the report from the spot, the testimony, to testify, to report
  \item the statement, the spokesperson, the press conference, the press release, the denial
  \item the edition, the caption, the scoop, censorship, to censor
  \item the live broadcast, the programme, to broadcast, the TV channel, the broadcaster
  \item the announcement, the statistic, to reveal, to spread, to go into in depth
\end{itemize}

\begin{tcolorbox}[breakable,colback=teal!4,colframe=teal!35!black,title={Before you move on}]
The eve? (\textbf{la vigilia}.) The next day? (\textbf{l'indomani}.) This morning? (\textbf{stamattina}.) Last night? (\textbf{stanotte}.) The whole morning? (\textbf{la mattinata}.) The news pages? (\textbf{la cronaca}.) The daily newspaper? (\textbf{il quotidiano}.) The weekly? (\textbf{il settimanale}.) The headline? (\textbf{il titolo}.) The article? (\textbf{l'articolo}.) The newsroom? (\textbf{la redazione}.) The news reporter? (\textbf{il cronista}.) The correspondent? (\textbf{l'inviato}.) The interview? (\textbf{l'intervista}.) To interview? (\textbf{intervistare}.) The inquiry? (\textbf{l'inchiesta}.) The report from the spot? (\textbf{il reportage}.) The testimony? (\textbf{la testimonianza}.) To testify? (\textbf{testimoniare}.) To report? (\textbf{riferire}.) The statement? (\textbf{la dichiarazione}.) The spokesperson? (\textbf{il portavoce}.) The press conference? (\textbf{la conferenza stampa}.) The press release? (\textbf{il comunicato}.) The denial? (\textbf{la smentita}.) The edition? (\textbf{l'edizione}.) The caption? (\textbf{la didascalia}.) The scoop? (\textbf{lo scoop}.) Censorship? (\textbf{la censura}.) To censor? (\textbf{censurare}.) The live broadcast? (\textbf{la diretta}.) The programme? (\textbf{la trasmissione}.) To broadcast? (\textbf{trasmettere}.) The TV channel? (\textbf{il canale}.) The broadcaster? (\textbf{l'emittente}.) The announcement? (\textbf{l'annuncio}.) The statistic? (\textbf{la statistica}.) To reveal? (\textbf{rivelare}.) To spread? (\textbf{diffondere}.) To go into in depth? (\textbf{approfondire}.)
\end{tcolorbox}
